\completa{1}{Arbués}{Qué estudiamos y qué vamos a responder}
\senalar{Señalar el título, el periodo y las dos cantidades inferiores.}
\decir{Buenos días. Somos Arbués Enrique Pérez Villanueva y Sergio Pezo. Presentamos el análisis de denuncias por desaparición de personas en el Perú, entre 2019 y 2024. Nuestra exposición sigue una misma idea: primero formulamos una pregunta que los datos pueden responder; luego mostramos las operaciones utilizadas y, finalmente, contrastamos la respuesta con la salida del programa.}
\decir{En la portada aparecen noventa y tres mil ochocientos treinta y un registros agregados y ciento veintiún mil cuatrocientas veintisiete personas contabilizadas. Son cantidades diferentes: un registro puede contener un conteo de varias personas. Por eso distinguimos entre contar filas y sumar CANTIDAD. Las veinticinco preguntas se trabajan con Scala, el lenguaje del código; Spark, el motor que ejecuta las operaciones; y PostgreSQL, que conserva una copia estructurada en una base relacional. A continuación ubicaremos cada parte dentro de la práctica.}

\completa{2}{Arbués}{Cómo se relaciona cada pregunta con la práctica}
\senalar{Recorrer todas las filas del mapa. Los números de la derecha son enlaces a fichas, no resultados numéricos.}
\decir{Este cuadro permite reconocer qué parte de la práctica resuelve cada grupo. Los ítems uno y dos presentan la fuente, los campos y las reglas de lectura. Los ítems tres a y tres b trabajan extremos y estadísticos mediante Scala y RDD. RDD significa colección distribuida resiliente: una colección que Spark divide en particiones y cuyas transformaciones puede reconstruir a partir de su origen.}
\decir{El punto tres c solicita tres consultas con resultados decimales y al menos tres MapReduce por consulta. En sus fichas identificaremos las agregaciones principales, las reducciones adicionales y el cálculo decimal. Así podremos señalar las operaciones que existen en los scripts y distinguirlas de una simple división. El punto cuatro corresponde al almacenamiento en PostgreSQL y al proceso ETL: extracción, transformación y carga.}
\decir{Después viene el punto cinco. En cinco a agrupamos y filtramos por dos o más columnas; en cinco b calculamos un promedio y un orden descendente. En cinco c mostramos cuatro formas de unir tablas: INNER, LEFT, RIGHT y FULL OUTER. En cinco d utilizamos funciones de Spark, disponibles mediante la importación functions punto guion bajo. En cinco e hacemos tres consultas con JOIN sobre vistas temporales; en cinco f ordenamos resultados combinando filtros; y en cinco g exportamos esas seis consultas a Excel.}
\decir{Los puntos seis y siete motivan la forma de presentar la solución: cada ficha contiene pregunta, entrada, operaciones, código y respuesta. Los enlaces permiten abrir las veinticinco consultas. El recorrido principal selecciona ejemplos; las demás también tienen su explicación en anexos. Antes de leer resultados, necesitamos precisar qué representa el dataset.}

\completa{3}{Arbués}{Fuente, alcance y unidad de análisis}
\senalar{Leer ambos bloques, la fórmula y el recuadro inferior.}
\decir{En el primer bloque leemos «Fuente y alcance». Los datos corresponden al Ministerio del Interior, abreviado MININTER, y se publican en la Plataforma Nacional de Datos Abiertos, o PNDA, dentro del rubro Seguridad y Justicia. Utilizamos el archivo local de enero de 2019 a diciembre de 2024: setenta y dos meses, noventa y tres mil ochocientas treinta y una filas, trece campos y mil quinientos seis ubigeos observados. Ubigeo es el código que identifica una ubicación administrativa. Las veintiséis unidades departamentales incluyen Lima Metropolitana y Región Lima por separado.}
\decir{El segundo bloque dice «Unidad de análisis». Una fila contiene un conteo para una combinación de año, mes, territorio y perfil. Dos filas del mismo distrito y mes pueden diferenciarse por sexo o rango de edad. No estamos leyendo una lista de personas con identificador individual. Por eso la fórmula dice: personas contabilizadas igual a la suma de CANTIDAD. Esa suma produce ciento veintiún mil cuatrocientas veintisiete. En cambio, count asterisco cuenta las filas: noventa y tres mil ochocientas treinta y una.}
\decir{Finalmente, el recuadro «Qué podemos interpretar» fija el alcance de las respuestas. Podemos describir cómo se distribuyen y cambian los conteos reportados. No podemos deduplicar personas, porque no hay identificadores individuales; tampoco establecer que desaparecieron juntas, explicar causas o calcular riesgo poblacional. Para una tasa de riesgo necesitaríamos, además, la población expuesta de cada grupo. Esta distinción acompañará todas las consultas. Ahora revisaremos cómo se leen las columnas.}

\completa{4}{Arbués}{Los trece campos y las reglas de lectura}
\senalar{Seguir todas las filas de la tabla y terminar en Limpieza y Derivadas.}
\decir{Para transformar esas preguntas en operaciones, revisamos los campos disponibles. AÑO y MES son enteros, números sin parte decimal; permiten agrupar por año, mes o trimestre. Cuando necesitamos una fecha, añadimos el día uno como convención: no es el día real de una desaparición. UBIGEO HECHO es texto, aunque contenga dígitos, para conservar los ceros iniciales de su código de seis posiciones. DPTO HECHO, PROV HECHO y DIST HECHO identifican departamento, provincia y distrito del hecho reportado.}
\decir{SEXO contiene F, M o NO ESPECIFICA. RANGO EDAD permite reconocer menores de dieciocho años mediante los intervalos de cero a cinco, seis a once y doce a diecisiete. EST CIVIL describe el estado civil; NACIONALIDAD permite separar las nacionalidades declaradas. Extranjero significa nacionalidad distinta de PERÚ y distinta de NO ESPECIFICA: un dato desconocido no se clasifica automáticamente como extranjero. NIVEL EDUCATIVO y OCUPACION describen otras características del perfil. Finalmente, CANTIDAD es un entero y una medida aditiva: podemos sumar sus valores. En la evidencia varía entre uno y dieciocho.}
\decir{La limpieza conserva los vacíos categóricos bajo la etiqueta NO ESPECIFICA. No imputamos esos atributos; imputar significaría asignarles un valor estimado. También generamos dos columnas derivadas. Región utiliza la clasificación por departamento definida en el código; no mide ruralidad. Periodo separa 2019, 2020 a 2021 y 2022 a 2024, con las etiquetas pre pandemia, pandemia y post pandemia. Son agrupaciones temporales de uno, dos y tres años.}
\decir{Al pie aparecen los nombres normalizados de Scala: anio, mes, ubigeo, departamento, provincia, distrito, sexo, rango edad, estado civil, nacionalidad, nivel educativo, ocupacion y cantidad. Los nombres con guiones bajos representan los mismos campos que acabamos de definir. El archivo esquema punto Scala centraliza su lectura. Veamos cómo ese mismo origen alimenta distintas formas de consulta.}

\completa{5}{Arbués}{Leer todas las ramas del flujo}
\senalar{Comenzar en CSV; seguir RDD, DataFrame y PostgreSQL. Las dos flechas que salen de DataFrame llegan a API y SQL.}
\decir{El cuadro inicial dice «CSV MININTER: trece campos, noventa y tres mil ochocientas treinta y una filas». CSV significa archivo de valores separados por comas. Desde él se abren tres ramas. Las flechas indican de dónde obtiene sus datos cada representación; no son pasos que deban ejecutarse todos en una sola cadena.}
\decir{La rama superior llega a «RDD: colección distribuida; parseo en Lc3 punto parsear». Parsear significa convertir una línea de texto en campos utilizables, como año entero y cantidad entera. La flecha siguiente conduce a «tres a a tres c: Scala y MapReduce; map, reduceByKey, reduce». Map transforma cada elemento; reduceByKey combina valores que comparten una clave; reduce combina todos los elementos de una colección en un resultado. MapReduce resume este patrón de transformar y combinar.}
\decir{En el centro leemos «DataFrame: tabla con esquema; limpieza y dimensiones». El esquema define nombres y tipos de columnas. Las dimensiones son catálogos descriptivos, como el que relaciona ubigeo y ubicación. La flecha horizontal llega a «cinco a a cinco d: API de DataFrame; groupBy, filter, join, funciones». API significa interfaz de programación: los métodos utilizados desde Scala. GroupBy agrupa, filter selecciona y join une por una condición.}
\decir{Desde DataFrame sale otra flecha hacia «cinco e a cinco f: vistas y spark punto sql; registrar vistas y consultar en SQL; cinco g: escribir resultados en Excel». Una vista temporal da un nombre consultable a un DataFrame dentro de la sesión de Spark. SQL es el lenguaje con el que expresamos qué datos queremos seleccionar, agrupar o unir. Esta rama utiliza Spark SQL; no consulta PostgreSQL.}
\decir{La rama inferior del CSV llega a «PostgreSQL dieciséis: extracción, transformación y carga, ETL». Esa flecha representa una carga independiente a la base. No hay una flecha desde PostgreSQL hacia las consultas Spark. Finalmente, el bloque inferior recuerda cómo leer los diagramas siguientes: Map crea claves y valores; Reduce los combina; join reúne agregados por clave y un map posterior calcula el porcentaje. A continuación veremos ese patrón con máximos y mínimos.}
